\documentclass[12pt]{article}
\usepackage[spanish]{babel}
\usepackage{geometry}
\geometry{a4paper, margin=2cm}
\usepackage{hyperref}
\usepackage{graphicx}
\usepackage{titlesec}
\usepackage{float}
\usepackage{enumitem}
\usepackage{fancyhdr}
\usepackage{setspace}
\usepackage{listings}       % Para códig
\usepackage{xcolor}         % Para colores en listings
\usepackage{booktabs}       % Para tablas profesionales
\usepackage{array}          % Para tablas mejoradas
\usepackage{multirow}       % Para celdas combinadas

\lstset{
	language=Python,
	basicstyle=\ttfamily\small,
	keywordstyle=\color{blue},
	commentstyle=\color{gray},
	stringstyle=\color{red},
	showstringspaces=false,
	breaklines=true,
	frame=single,
	numbers=left,
	numberstyle=\tiny\color{gray}
}



\pagestyle{fancy}
\fancyhf{}
\rhead{ML}
\lhead{Informe Técnico}
\rfoot{\thepage}
\setlength{\headheight}{14pt}

\title{%
	\vspace{1cm}
	\Huge \textbf{Etiquetado con modelo no supervisado}\\[6pt]
	\Large \textbf{Diseño e Implementación de una Red Social Distribuida}\\[12pt]
}
\author{\textit{Victor Hugo Pacheco Fonseca C411}\\
	\textit{Jose Agustin del Toro C412}\\
	Asignatura: Machine Learning}
\date{2025}

\begin{document}
	
	\maketitle
	\tableofcontents
	\newpage
	
	%------------------------------------------------------------
	\section{Resumen}
	Se implementó un sistema híbrido de detección de anomalías en logs web Apache combinando técnicas de machine learning no supervisado (Isolation Forest) con un sistema de reglas heurísticas. El proceso incluyó etiquetado semiautomático de datos, ingeniería de características, optimización de hiperparámetros, validación manual de resultados y enriquecimiento iterativo del sistema de reglas.
	
	\section{Metodología de Etiquetado de Datos}
	\subsection{Dataset Inicial}
	\begin{itemize}
		\item \textbf{Origen:} Logs de acceso Apache de servidores web productivos
		\item \textbf{Período:} Noviembre 2025
		\item \textbf{Tamaño}: \textasciitilde 1,000,000 registros
		\item \textbf{Estado:} Aproximadamente 50\% pre-etiquetado mediante sistema de reglas existente
	\end{itemize}

	\section{Ingeniería de Características para Isolation Forest}
	\begin{itemize}
		\item \textbf{A. Características Numéricas (StandardScaler):}
		\begin{itemize}
			\item \textbf{Frecuencias y ocurrencias:}
			\begin{itemize}
				\item \texttt{uri\_occurrence}: Frecuencia absoluta de cada URL
				\item \texttt{user\_agent\_occurrence}: Frecuencia de cada User-Agent
				\item \texttt{ip\_occurrence}: Frecuencia por dirección IP
				\item \texttt{ioc\_occurrence}: Conteo de Indicators of Compromise
			\end{itemize}
			
			\item \textbf{Características de longitud:}
			\begin{itemize}
				\item \texttt{uri\_length}: Longitud del string URI
				\item \texttt{user\_agent\_length}: Longitud del User-Agent
				\item \texttt{uri\_decoded\_length}: Longitud después de URL decode
			\end{itemize}
			
	
			
			\item \textbf{Estadísticas de respuesta:}
			\begin{itemize}
				\item \texttt{response\_size}: Tamaño de respuesta en byte
			\end{itemize}
		\end{itemize}
		
		\item \textbf{B. Características Categóricas (OneHotEncoding):}
		\begin{itemize}
			\item \texttt{method\_GET}, \texttt{method\_POST}, \texttt{method\_HEAD}
			\item \texttt{protocol\_HTTP\_1.1}, \texttt{protocol\_HTTP\_2.0}
			\item \texttt{status\_2xx}, \texttt{status\_4xx}, \texttt{status\_5xx}
		\end{itemize}
	\end{itemize}
	



	\section{Implementación y Optimización de Isolation Forest}
	
	\subsection{Configuración del Modelo}
	
	\begin{verbatim}
		IsolationForest(
		n_estimators=100,
		max_samples=0.25,      # 25% para reducir swamping/masking
		max_features=0.5,      # 50% de características por árbol
		bootstrap=False,
		n_jobs=-1,            # Paralelización completa
		random_state=42       # Reproducibilidad
		)
	\end{verbatim}
	
	\subsection{Experimentación con Parámetro Contamination}
	
	Se evaluaron 10 valores de \texttt{contamination} en un dataset de 576,767 registros con 84,530 anomalías reales (14.65\%):
	
	\begin{table}[htbp]
		\centering
		\caption{Resultados de experimentación con diferentes valores de contamination}
		\label{tab:contamination-results}
		\begin{tabular}{|c|c|c|c|c|c|c|c|}
			\hline
			\textbf{Contamination} & \textbf{Accuracy} & \textbf{Precision} & \textbf{Recall} & \textbf{F1-Score} & \textbf{TP} & \textbf{FP} & \textbf{FN} \\
			\hline
			0.001 (0.1\%) & 0.8544 & 0.9705 & 0.0066 & 0.0132 & 560 & 17 & 83,970 \\
			\hline
			0.005 (0.5\%) & 0.8548 & 0.6397 & 0.0218 & 0.0422 & 1,845 & 1,039 & 82,685 \\
			\hline
			0.01 (1\%) & 0.8568 & 0.6657 & 0.0454 & 0.0850 & 3,839 & 1,928 & 80,691 \\
			\hline
			0.02 (2\%) & 0.8588 & 0.6351 & 0.0867 & 0.1525 & 7,327 & 4,209 & 77,203 \\
			\hline
			\textbf{0.03 (3\%)} & \textbf{0.8564} & \textbf{0.5486} & \textbf{0.1123} & \textbf{0.1864} & \textbf{9,492} & \textbf{7,811} & \textbf{75,038} \\
			\hline
			0.05 (5\%) & 0.8617 & 0.5823 & 0.1986 & 0.2962 & 16,786 & 12,040 & 67,744 \\
			\hline
			0.075 (7.5\%) & 0.8474 & 0.4599 & 0.2353 & 0.3113 & 19,893 & 23,364 & 64,637 \\
			\hline
			0.1 (10\%) & 0.8512 & 0.4879 & 0.3079 & 0.3775 & 26,026 & 27,319 & 58,504 \\
			\hline
			0.15 (15\%) & 0.8718 & 0.5612 & 0.5743 & 0.5677 & 48,549 & 37,965 & 35,981 \\
			\hline
			0.2 (20\%) & 0.8548 & 0.5033 & 0.6865 & 0.5808 & 58,026 & 57,256 & 26,504 \\
			\hline
		\end{tabular}
	\end{table}
	
	\subsection{Análisis de los Resultados}
	
	\begin{itemize}
		\item \textbf{Alto Precision, Bajo Recall (0.001-0.01):} 
		\begin{itemize}
			\item Contamination 0.001: Precision 97.05\% pero Recall solo 0.66\%
			\item Detecta pocas anomalías pero con alta confianza
			\item Ideal para entornos donde falsos positivos son críticos
		\end{itemize}
		
		\item \textbf{Zona de Balance (0.03-0.1):}
		\begin{itemize}
			\item Contamination 0.03: F1-Score 0.1864, Precision 54.86\%, Recall 11.23\%
			\item Contamination 0.05: Mejor equilibrio con F1-Score 0.2962
			\item Recall mejora significativamente mientras Precision disminuye gradualmente
		\end{itemize}
		
		\item \textbf{Alto Recall, Bajo Precision (0.15-0.2):}
		\begin{itemize}
			\item Contamination 0.15: Recall 57.43\%, Precision 56.12\%, F1-Score 0.5677
			\item Contamination 0.2: Recall 68.65\% (más alto) pero Precision baja a 50.33\%
			\item Detecta más anomalías pero genera más falsos positivos
		\end{itemize}
	\end{itemize}
	
	\section{Verificación Manual de Resultados}
	 Una vez obtenido el puntaje de anomalía del modelo, se ordenaron los registros según su \textit{anomaly score}. Se revisaron manualmente los casos con valores más bajos (más anómalos) para determinar:
	  \begin{itemize}
	  	 \item Si realmente correspondían a ataques. 
	  	 \item Si existía un patrón repetido que justificara crear una regla específica. 
	  \end{itemize} 
	  Este proceso permitió: \begin{itemize}
	  	 \item Confirmar ataques reales detectados por el modelo.
	  	  \item Identificar falsos positivos. 
	  	  \item Descubrir nuevas familias de ataques no contempladas inicialmente. 
	  	  \end{itemize} 
	  	  \section{Creación de Reglas de Detección} 
	  	  Cuando se identificaba un ataque real con un patrón claro, se diseñaba una regla para detectarlo automáticamente. Ejemplos de patrones incorporados:
	  	   \begin{itemize}
	  	   	 \item Inyecciones SQL con operadores como \texttt{CAST}, \texttt{CASE WHEN}, \texttt{||}, \texttt{::text}.
	  	   	  \item URLs con comentarios ofuscados (\texttt{/* ... */}). 
	  	   	  \item Terminaciones típicas de SQLi (\texttt{--+}). \item Bots maliciosos con \textit{User-Agent}
	  	   	   falsificado.
	  	   	    \end{itemize} 
	  	   	    Cada regla se añadía al sistema para mejorar la detección futura y reducir la dependencia del análisis manual.
	

\section{Conclusiones} El proceso permitió:
 \begin{itemize}
 	 \item Mejorar la calidad del etiquetado del dataset.
 	  \item Detectar ataques complejos que no estaban cubiertos por reglas previas. 
 	  \item Ajustar el modelo de \textit{Isolation Forest} para maximizar su utilidad. 
 	  \item Crear reglas específicas que fortalecen la seguridad del sistema. 
 	  \end{itemize}
 	   El enfoque híbrido (modelo + verificación manual + reglas) demostró ser efectivo para la detección progresiva de nuevas amenazas.


\end{document}